\documentclass{article}
\usepackage[russian]{babel}
\usepackage{paracol}
\usepackage{tikz}

\begin{document}

\begin{paracol}{2}
\noindent 14 \\
\begin{flushright}

\vspace{12cm}
\begin{tikzpicture}[scale=0.8]
    \draw[orange, line width=2pt] (0,0) -- (2,0);
    \draw[blue, line width=2pt] (0,0) -- (1.5,1.5);
    \fill[yellow] (0,0) -- (0.8,0) arc (0:45:0.8);
\end{tikzpicture} \\

\vspace{2cm}

\begin{tikzpicture}[scale=0.8]
    \draw[black, line width=2pt] (-1.5,0) -- (1.5,0);
    \draw[black, line width=2pt] (0,0) -- (0,1);
    \draw[black, line width=2pt] (0.8,0) arc (0:180:0.8);
\end{tikzpicture} \\
\end{flushright}

\switchcolumn

\begin{center}
\textbf{1} \\
Точка есть то, что не имеет частей. \\[1em]

\textbf{2} \\
Линия — это длина без ширины. \\[1em]

\textbf{3} \\
Концы линии — точки. \\[1em]

\textbf{4} \\
Прямая линия есть та, которая равно расположена по отношению к точкам на ней. \\[1em]

\textbf{5} \\
Поверхность есть то, что имеет только длину и ширину. \\[1em]

\textbf{6} \\
Концы поверхности — линии. \\[1em]

\textbf{7} \\
Плоская поверхность есть та, которая равно расположена по отношению к прямым на ней. \\[1em]

\textbf{8} \\
Плоский угол есть наклонение друг к другу двух линий, в плоскости встречающихся друг с другом, но не направленных одинаково. \\[1em]

\textbf{9} \\  
Прямолинейный угол есть наклонение друг к другу двух пересекающихся и не совпадающих прямых линий. \\[1em]

\textbf{10} \\
Когда прямая восстановленная на другой прямой образует равные смежные углы, каждый из этих углов называется \textit{прямым углом}, а каждая из этих прямых называется \textit{перпендикуляром} к другой.


\end{center}

\end{paracol}

\end{document}
